\section{Introducción}

Durante el desarrollo de sus actividades diarias, el personal que maneja fuentes de radiación debe estar consciente de que, si bien las radiaciones ionizantes son herramientas útiles en la industria, la medicina y la investigación, también son potencialmente peligrosas si no se manejan de manera adecuada.

Por lo tanto, es de suma importancia conocer y respetar una serie de normas y protocolos que garanticen la seguridad radiológica y protejan tanto a los trabajadores como al público en general. La implementación de principios básicos de seguridad y protección radiológica, como el tiempo, la distancia y el blindaje, es fundamental para minimizar la exposición innecesaria a las radiaciones.

Este manual está diseñado para proporcionar la información esencial y las instrucciones operativas necesarias para asegurar que todas las actividades que involucran el uso de radiación se realicen de manera segura. Además, se promueve una cultura de seguridad mediante la cual cada individuo se haga responsable de su propia seguridad y la de los demás.

\subsection{Principios Básicos de Protección Radiológica}

Los principios básicos de protección radiológica incluyen:

\begin{itemize}
  \item \textbf{Justificación:} Cualquier práctica que involucre una exposición a la radiación debe ser justificada, en el sentido de que los beneficios para la sociedad sean mayores que el riesgo asociado.
  \item \textbf{Optimización:} Las exposiciones deben ser mantenidas tan bajas como sea razonablemente posible, tomando en consideración factores económicos y sociales. Este principio es conocido como ALARA (As Low As Reasonably Achievable).
  \item \textbf{Limitación de dosis:} Las dosis individuales no deben exceder los límites establecidos en las normativas nacionales e internacionales.
\end{itemize}

Estos principios no solo se aplican a los profesionales que trabajan directamente con fuentes de radiación, sino también a la protección del público general y del medio ambiente.